\documentclass[11pt]{article}
\usepackage[margin=1in]{geometry}
\usepackage{amsmath,amssymb,amsthm,mathtools}
\usepackage[colorlinks=true,linkcolor=blue,citecolor=blue,urlcolor=blue]{hyperref}

\newtheorem{theorem}{Theorem}
\renewcommand{\thetheorem}{\Alph{theorem}}
\newtheorem{proposition}{Proposition}[section]
\newtheorem{lemma}[proposition]{Lemma}
\newtheorem{corollary}[proposition]{Corollary}
\theoremstyle{remark}
\newtheorem{remark}[proposition]{Remark}

\newcommand{\eps}{\varepsilon}
\newcommand{\cA}{\mathcal{A}}
\newcommand{\cB}{\mathcal{B}}
\newcommand{\cH}{\mathcal{H}}
\newcommand{\cK}{\mathcal{K}}
\newcommand{\cL}{\mathcal{L}}
\newcommand{\cE}{\mathcal{E}}
\newcommand{\Bo}{\mathbf{B}}
\newcommand{\cb}{\mathrm{cb}}
\newcommand{\Img}{\operatorname{Im}}
\newcommand{\wt}{\widetilde}
\newcommand{\Ma}[1]{\mathbf{M}_{#1}}

\title{Kitaev's $\eps$-$C^*$ rigidity beyond finite dimensions:\\ what extends, what is Lean-checked, and what is open}
\author{Joydeep Naskar\\[2pt]
{\small Beijing Institute of Mathematical Sciences and Applications (BIMSA)}\\
{\small Huaibei Town, Huairou District, Beijing 101408, China}\\
{\small\href{mailto:joydeepnaskar@bimsa.cn}{\texttt{joydeepnaskar@bimsa.cn}}}
\and
Claude\\[2pt]
{\small Anthropic}}
\date{October 18, 2026}

\begin{document}
\maketitle

\begin{abstract}
Kitaev showed that a finite-dimensional $\eps$-$C^*$ algebra, which satisfies the $C^*$ axioms up to errors of order $\eps$, is $O(\eps)$-isomorphic to a genuine $C^*$ algebra, and that an almost-idempotent quantum channel on a finite-dimensional Hilbert space approximately factors through such an algebra. We study which of these statements survive in infinite dimensions. When the Hilbert space is infinite-dimensional but the almost-invariant algebra is finite-dimensional, as for codes with a finite-dimensional logical system, both results hold with the same universal constants, since the two Choi representations in Kitaev's proof can be replaced by Stinespring dilations. In the course of the proof we correct two points of the published argument, namely the diagonal of a block algebra with two or more blocks and the exact case of an idempotent channel. The results also extend to $\ell^\infty$-direct sums of finite-dimensional $\eps$-$C^*$ algebras, because Kitaev's constants are uniform. For a general infinite-dimensional $\eps$-$C^*$ algebra the question is open. We identify where Kitaev's proof uses finite dimension, show that there is no local obstruction, and point to degree-three bounded cohomology of non-amenable groups as the likely source of a counterexample. Both positive results are verified in Lean~4 with Mathlib, conditional on the cited results of Kitaev's paper and on standard facts about completely positive maps.
\end{abstract}

\noindent
All equation, theorem and section numbers of the form (12.37), Theorem~2.3, \S12 refer to Kitaev, \emph{Almost-idempotent quantum channels and approximate $C^*$-algebras}, arXiv:2405.02434v2.

\section{The question, and the short answer}

Kitaev's Theorem~2.3 says that there are universal constants $\eps_0, C$ such that every \emph{finite-dimensional} $\eps$-$C^*$ algebra $\cA$ with $\eps<\eps_0$ admits a linear bijection $v\colon\cB\to\cA$ from a genuine $C^*$ algebra $\cB$ which is a $C\eps$-isomorphism: $v(X^\dagger)=v(X)^\dagger$, $\|v(I)-I\|\le C\eps$, $\|v(XY)-v(X)v(Y)\|\le C\eps\|X\|\|Y\|$ and $(1-C\eps)\|X\|\le\|v(X)\|\le(1+C\eps)\|X\|$. The question is whether this survives in infinite dimensions. ``Infinite dimensions'' can mean two different things, and the answers differ.

\begin{itemize}
\item[(a)] \textbf{The physical Hilbert space $\cH$ is infinite-dimensional, but the almost-invariant algebra $\cA=\Img\wt\Phi$ is finite-dimensional.} This is the situation of bosonic or field-theoretic codes with a finite-dimensional logical system. Here the answer is \emph{yes}: both the bijection $v$ and Kitaev's UCP factorization $\Phi\approx\Delta\Upsilon$ (Theorem~12.3) hold, with the same universal constants and no assumption on $\dim\cH$ (Theorem~\ref{thm:A}). Kitaev states Theorem~12.3 only for $\dim\cH<\infty$; the proof uses finite dimensionality of $\cH$ only through two Choi representations, and both can be replaced by Stinespring dilations. Theorem~\ref{thm:A} is machine-checked in Lean, conditional on the cited results of Kitaev's paper and three textbook facts about completely positive maps (\S\ref{sec:lean}). The formalization exposed two points in the published argument that need repair, and both are fixed below.
\item[(b)] \textbf{The $\eps$-$C^*$ algebra $\cA$ itself is infinite-dimensional.} If $\cA$ is an $\ell^\infty$-direct sum of finite-dimensional pieces (infinitely many finite sectors) the answer is \emph{yes}, because Kitaev's constants are uniform (Theorem~\ref{thm:B}). We have verified this reduction in Lean~4 with Mathlib, with Kitaev's finite-dimensional theorem as an explicit hypothesis (\S\ref{sec:lean}). For a general infinite-dimensional $\cA$ the question is \emph{open}. We can neither prove nor refute it, and we found no claim either way in the literature (Kitaev; Ranard--Walter--Witteveen, arXiv:2603.08702, who use the finite-dimensional theorem; Alekseev--Farah--Thom, arXiv:2606.03757, who remark that $\eps$-$C^*$ algebras are ``subtly different'' from Ulam stability). Section~\ref{sec:open} explains where Kitaev's proof needs $\dim\cA<\infty$. It also shows that no counterexample can be detected on finitely many elements, and identifies the most plausible place one could come from: degree-3 bounded cohomology of non-amenable groups.
\end{itemize}

\section{Infinite-dimensional $\cH$, finite-dimensional logical algebra}

Throughout, $\Phi\colon\Bo(\cH)\to\Bo(\cH)$ is UCP, $\|\Phi^2-\Phi\|_\cb\le\eta$, $\wt\Phi=\theta(2\Phi-1)$ is defined by (12.2), $\cA=\Img\wt\Phi$, and $X\star Y=\wt\Phi(XY)$.

\begin{lemma}\label{lem:rank}
If $\eta<1/4$ then $\wt\Phi=\Phi\Psi$ for a bounded map $\Psi$ commuting with $\Phi$. Hence $\Img\wt\Phi\subseteq\Img\Phi$ and $\operatorname{rank}\wt\Phi\le\operatorname{rank}\Phi$.
\end{lemma}

\begin{proof}
Let $N=4(\Phi-\Phi^2)=4\Phi(1-\Phi)$, so $\|N\|_\cb\le 4\eta<1$, and $(1-x)^{-1/2}=\sum_{k\ge0}c_kx^k$ with $c_0=1$. Then (12.2) reads
\[
\wt\Phi=\tfrac12\Bigl(1+(2\Phi-1)\sum_{k\ge0}c_kN^k\Bigr)
=\Phi+(\Phi-\tfrac12)\sum_{k\ge1}c_kN^k
=\Phi\Bigl[1+4(\Phi-\tfrac12)(1-\Phi)\sum_{k\ge1}c_kN^{k-1}\Bigr],
\]
using $N^k=\Phi\cdot4(1-\Phi)N^{k-1}$. The bracket converges absolutely and commutes with $\Phi$.
\end{proof}

\begin{theorem}\label{thm:A}
There are universal constants $\eta_0,C>0$ with the following property. Let $\cH$ be a nonzero Hilbert space of arbitrary dimension and $\Phi\colon\Bo(\cH)\to\Bo(\cH)$ a UCP map with $\|\Phi^2-\Phi\|_\cb\le\eta\le\eta_0$ such that $\wt\Phi$ has finite rank (by Lemma~\ref{lem:rank} it suffices that $\Phi$ has finite rank). Then
\begin{enumerate}
\item[(i)] there are a finite-dimensional $C^*$ algebra $\cB$ and an extended $C\eta$-isomorphism $v\colon\cB\to\cA$; in particular $v$ is a linear bijection;
\item[(ii)] there are UCP maps $\Delta\colon\cB\to\Bo(\cH)$ and $\Upsilon\colon\Bo(\cH)\to\cB$ with $\|\Delta\Upsilon-\Phi\|_\cb\le C\eta$ and $\|\Upsilon_n(\Delta_n(X)\Delta_n(Y))-XY\|\le C\eta\|X\|\|Y\|$ for all $n$ and $X,Y\in\Ma{n}\otimes\cB$; in particular $\|\Upsilon\Delta-1_\cB\|_\cb\le C\eta$.
\end{enumerate}
\end{theorem}

\begin{proof}
(i) Kitaev proves Theorem~12.1 for arbitrary $\cH$; the ``Stinespring stack'' of \S12.1 is there for exactly this purpose. So $\cA$, with the matrix norms inherited from $\Ma{n}(\Bo(\cH))$, is an extended $O(\eta)$-$C^*$ algebra. It is finite-dimensional by hypothesis. Theorem~10.5 is a statement about abstract finite-dimensional extended $\eps$-$C^*$ algebras and makes no reference to $\cH$, so it gives $\cB$ and $v$.

(ii) Kitaev's proof of Theorem~12.3 consists of the outline (12.25)--(12.27), the construction of $\Delta'$ and $\Delta$, the properties (12.33)--(12.36), Lemma~12.4, and the construction of $\Upsilon'$ and $\Upsilon$. The outline, $\Delta$, and (12.33)--(12.36) use only the diagonal of the finite-dimensional algebra $\cB$, the bounds of Theorem~12.1 (including (12.6) applied to $\Phi_n$), and normalization by $\Delta'(I)^{-1/2}$. None of these uses $\dim\cH<\infty$.

One correction is needed in this part. Kitaev represents the diagonal of each block $\Bo(\cL_j)$ by a unitary 1-design $D_j=\sum_sp_{js}U_{js}^\dagger\otimes U_{js}$ satisfying (12.29), and takes as the diagonal of $\cB$ the product $D=\sum_sp_sU_s^\dagger\otimes U_s$, with $p_s=\prod_jp_{js_j}$ and $U_s=\bigoplus_jU_{js_j}$. For two or more blocks this $D$ does not satisfy $XD=DX$. For $\cB=\mathbb{C}\oplus\mathbb{C}$ with trivial designs, $D=1\otimes1$, and $X=(1,0)$ gives $XD\neq DX$. The cross-block terms survive because they carry the first moments $\sum_sp_{js}U_{js}$, which need not vanish. The identity $XD=DX$ is what makes $\Delta'$ completely positive. It is restored by attaching an independent uniform sign to each block, $U_{s,\sigma}=\bigoplus_j\sigma_jU_{js_j}$ with $\sigma\in\{\pm1\}^m$ uniformly distributed. Every cross-block term then averages to zero, each diagonal block reproduces (12.29), and nothing else in the argument changes.

Finite dimensionality of $\cH$ enters only through two Choi representations, which we replace as follows.

\emph{(a) The representation (12.37) of $\Delta$.} Write $\cB=\bigoplus_{j=1}^m\Bo(\cL_j)$, $\dim\cL_j<\infty$. By Stinespring, $\Delta(X)=W^\dagger\rho(X)W$ with $\rho$ a unital $*$-representation of $\cB$ on some $\cK$ and $W$ an isometry. Every unital representation of $\bigoplus_j\Bo(\cL_j)$ is unitarily equivalent to $X\mapsto\bigoplus_j X_j\otimes1_{\cE_j}$ on $\bigoplus_j\cL_j\otimes\cE_j$, for Hilbert spaces $\cE_j$ of arbitrary dimension (possibly $0$). Writing $W=(W_1,\dots,W_m)$ gives exactly (12.37), with $\sum_jW_j^\dagger W_j=W^\dagger W=1_\cH$.

\emph{(b) Lemma~12.4.} The map $a\otimes b\mapsto(a\otimes1)W_jW_j^\dagger(b\otimes1)$ is linear on the algebraic tensor product $\Bo(\cL_j)\otimes\Bo(\cL_j)$. Applying it to (12.29) shows that $R_j$ in (12.39) commutes with $\Bo(\cL_j)\otimes1_{\cE_j}$. The commutant of $\Bo(\cL_j)\otimes1$ in $\Bo(\cL_j\otimes\cE_j)$ is $1\otimes\Bo(\cE_j)$ whatever the dimension of $\cE_j$, so $R_j=1\otimes C_j$. The bounds $1-O(\eta)\le\|C_j\|\le1$ use only (12.33) and (12.35). The one change is that, when $\dim\cE_j=\infty$, $\|C_j\|$ need not be attained. For $\eta>0$ we choose a unit vector $\xi_j$ with $\|C_j\xi_j\|\ge\|C_j\|-\eta\ge1-O(\eta)$, which changes only the constant. For $\eta=0$ this choice is not available, while the theorem then asks for an exact factorization; this case is treated separately in (d).

\emph{(c) The Choi representation $\Phi(X)=V^\dagger(X\otimes1_\mathcal{F})V$.} This exists only for normal $\Phi$. Instead take a Stinespring dilation $\Phi(X)=V^\dagger\pi(X)V$, with $\pi\colon\Bo(\cH)\to\Bo(\cK')$ a unital $*$-representation and $V$ an isometry, and define
\[
L_j=\sum_sp_{js}\,\pi\bigl(\Delta(U_{js}^\dagger)\bigr)\,V\,W_j^\dagger(U_{js}\otimes\xi_j)\colon\cL_j\to\cK',\qquad
\Upsilon'_j(X)=L_j^\dagger\,\pi\bigl(\Phi(X)\bigr)\,L_j .
\]
Each $\Upsilon'_j$ is completely positive and $\|L_j\|\le1$. Since $\pi(\Delta(U^\dagger))^\dagger=\pi(\Delta(U))$ and
\[
V^\dagger\pi(Z_1)\pi(Z_2)\pi(Z_3)V=\Phi(Z_1Z_2Z_3),
\]
we get
\[
\Upsilon'_j(\Delta(Y))=\sum_{r,s}p_{jr}p_{js}\,(U_{jr}^\dagger\otimes\xi_j^\dagger)\,W_j\,
\Phi\bigl(\Delta(U_{jr})\,\Phi(\Delta(Y))\,\Delta(U_{js}^\dagger)\bigr)\,W_j^\dagger(U_{js}\otimes\xi_j).
\]
This is the first line of Kitaev's computation of $\Upsilon'_j(\Delta(Y))$. The subsequent lines use (12.34), (12.36), (12.37), (12.39) and Lemma~12.4, all of which hold by (a) and (b). At matrix level $n$, the dilation $(1_n\otimes V,\,1_n\otimes\pi)$ of $\Phi_n$ gives $(\Upsilon'_j)_n(X)=(1\otimes L_j)^\dagger(1\otimes\pi)(\Phi_n(X))(1\otimes L_j)$, so the same computation applies. Likewise $\Upsilon'_j\Phi-\Upsilon'_j=L_j^\dagger\pi\bigl((\Phi^2-\Phi)(\cdot)\bigr)L_j$ has cb-norm at most $\eta$. The chain $\Upsilon'\approx\Upsilon'\Phi\approx\Upsilon'\wt\Phi=\Upsilon'\wt\Delta\wt\Upsilon\approx\Upsilon'\Delta\wt\Upsilon\approx\wt\Upsilon$ and the normalization $\Upsilon=\Upsilon'(1)^{-1/2}\Upsilon'(\cdot)\Upsilon'(1)^{-1/2}$ are dimension-free.

\emph{(d) The case $\eta=0$.} Then $\Phi^2=\Phi$, so $\wt\Phi=\Phi$, and Theorem~10.5 gives $\delta=0$. So $v$ is a unital complete isometry with $\Phi(v(X)v(Y))=v(XY)$, hence $\Delta'=v$ and $\Delta=v$. Take $\Upsilon=\wt\Upsilon=v^{-1}\Phi$. It is unital, and it is completely positive. If $X\ge0$ in $\Ma{n}(\Bo(\cH))$, then $Y=\Upsilon_n(X)$ is self-adjoint and $v_n(Y)=\Phi_n(X)\ge0$. With $t=\|Y\|=\|v_n(Y)\|$ we get $\|t-Y\|=\|t-v_n(Y)\|\le t$, and a self-adjoint element with $\|t-Y\|\le t$ is positive. Finally $\Delta\Upsilon=\Phi$ and $\Upsilon_n(\Delta_n(X)\Delta_n(Y))=v_n^{-1}\bigl(\Phi_n(v_n(X)v_n(Y))\bigr)=XY$ exactly.
\end{proof}

\begin{remark}
A simpler choice of $\xi_j$ also works, although it is not part of the Lean proof. The computation in (c) gives $\Upsilon'\Delta\approx D$ rather than $\approx1$, where $D=(\|C_j\xi_j\|^2)_j$, so the chain yields $\Upsilon'=D\wt\Upsilon+O(\eta)$. Since $D$ is central in $\cB$, $\Upsilon'(1)^{-1/2}=D^{-1/2}+O(\eta)$, and the normalization divides $D$ out. So any $\xi_j$ with $\|C_j\xi_j\|\ge1/2$ suffices, uniformly in $\eta\ge0$, and no approximate maximizer is needed.
\end{remark}

\begin{remark}
The rank hypothesis is about the spectrum of $\Phi$ near $1$, not about $\cH$. Since $\|\Phi^2-\Phi\|_\cb<1/4$, $\sigma(\Phi)$ splits into a part near $0$ and a part near $1$, and $\wt\Phi$ is the Riesz projection onto the latter. So $\wt\Phi$ has finite rank exactly when that part consists of finitely many eigenvalues of finite algebraic multiplicity. This holds, for instance, whenever $\Phi$ is a compact operator on the Banach space $\Bo(\cH)$. The case $\Phi=1$ (where $\cA=\Bo(\cH)$ is already a $C^*$ algebra) shows that the hypothesis is not automatic.
\end{remark}

\begin{corollary}[codes through a finite-dimensional system]\label{cor:codes}
Let $\mathcal{C}\subseteq\cH$ be a finite-dimensional subspace and $T$ a channel on trace-class operators whose outputs are supported in $\mathcal{C}$, for example noise followed by a recovery map into a code space. If $\Phi=T^*$ satisfies $\|\Phi^2-\Phi\|_\cb\le\eta\le\eta_0$, then (i) and (ii) of Theorem~\ref{thm:A} hold.
\end{corollary}

\begin{proof}
Since $T(\rho)=P_\mathcal{C}T(\rho)P_\mathcal{C}$, we have $\Phi(X)=\Phi(P_\mathcal{C}XP_\mathcal{C})$, so $\operatorname{rank}\Phi\le(\dim\mathcal{C})^2$. Apply Lemma~\ref{lem:rank} and Theorem~\ref{thm:A}.
\end{proof}

\section{Infinitely many finite sectors}

\begin{proposition}\label{prop:sum}
Let $I\neq\emptyset$. For each $i\in I$ let $\cA_i$ be an $\eps$-$C^*$ algebra, $\cB_i$ a $C^*$ algebra, and $v_i\colon\cB_i\to\cA_i$ a $\delta$-isomorphism, with the same $\eps$ and the same $\delta<1$ for all $i$. Let $\cA=\bigoplus^\infty_i\cA_i$ and $\cB=\bigoplus^\infty_i\cB_i$ be the bounded families with supremum norm and coordinatewise operations. Then $\cA$ is an $\eps$-$C^*$ algebra, $\cB$ is a $C^*$ algebra, and $v=(v_i)_i\colon\cB\to\cA$ is a $\delta$-isomorphism. The same holds for extended structures.
\end{proposition}

\begin{proof}
Every axiom (2.1)--(2.6) is an inequality between quantities computed coordinatewise. It holds in each coordinate with the same $\eps$, so taking suprema gives it in $\cA$. For instance $(1-\eps)\|a_i\|^2\le\|a_i^\dagger a_i\|\le\|a^\dagger a\|$ for every $i$, hence $(1-\eps)\|a\|^2\le\|a^\dagger a\|$. The bound $|\|I\|-1|\le\eps$ is where $I\ne\emptyset$ is needed: $\|I\|=\sup_i\|I_i\|\in[1-\eps,1+\eps]$. Next, $\|v(b)\|=\sup_i\|v_i(b_i)\|$ lies in $[(1-\delta)\|b\|,(1+\delta)\|b\|]$, so $v$ is well defined and bounded. The inverse $(v_i^{-1})_i$ has norm at most $(1-\delta)^{-1}$, so $v$ is bijective. The unit and multiplicativity defects are coordinatewise at most $\delta$ and $\delta\|b_i\|\|c_i\|\le\delta\|b\|\|c\|$. For extended structures use $\Ma{n}(\bigoplus^\infty\cA_i)=\bigoplus^\infty\Ma{n}(\cA_i)$ isometrically.
\end{proof}

\begin{theorem}\label{thm:B}
Let $\eps_0,C$ be the constants of Kitaev's Theorem~2.3 (or 10.5). Every $\ell^\infty$-direct sum of finite-dimensional (extended) $\eps$-$C^*$ algebras with $\eps<\eps_0$ and $C\eps<1$ is (extended) $C\eps$-isomorphic to a $C^*$ algebra. Consequently, if $\cH=\bigoplus_k\cH_k$ and $\Phi(X)=\bigoplus_k\Phi_k(P_kXP_k)$, where each $\Phi_k$ is UCP on $\Bo(\cH_k)$ with $\|\Phi_k^2-\Phi_k\|_\cb\le\eta\le\eta_0$ and $\wt\Phi_k$ of finite rank, then $\|\Phi^2-\Phi\|_\cb\le\eta$, and (i) and (ii) of Theorem~\ref{thm:A} hold with the infinite-dimensional algebra $\cB=\bigoplus^\infty_k\cB_k$.
\end{theorem}

\begin{proof}
The first claim follows from Proposition~\ref{prop:sum} by choosing, for each $i$, Kitaev's $\cB_i$ and $v_i$. For the channel statement, $\Phi^m(X)=\bigoplus_k\Phi_k^m(P_kXP_k)$ for $m\ge1$, and by the proof of Lemma~\ref{lem:rank}, $\wt\Phi=\Phi\Psi$ is a norm-convergent series in $\Phi$ without constant term. Hence $\wt\Phi(X)=\bigoplus_k\wt\Phi_k(P_kXP_k)$ and $\cA=\bigoplus^\infty_k\Img\wt\Phi_k$ with coordinatewise Choi--Effros product. Put $\Delta=\bigoplus_k\Delta_k$ and $\Upsilon(X)=(\Upsilon_k(P_kXP_k))_k$, both UCP. All the required bounds are suprema of the bounds of Theorem~\ref{thm:A} for the $\Phi_k$.
\end{proof}

\begin{remark}
The same argument shows that ultraproducts $\prod_\mathcal{U}\cA_i$ of finite-dimensional $\eps$-$C^*$ algebras are $C\eps$-isomorphic to $\prod_\mathcal{U}\cB_i$: the maps $v_i$ and $v_i^{-1}$ are uniformly bounded, so $\prod_\mathcal{U}v_i$ is a well-defined bijection, and the defects pass to the limit. Thus the finite-dimensional theorem with \emph{uniform} constants already covers every infinite-dimensional algebra that is assembled from finite-dimensional ones in a uniformly bounded way.
\end{remark}

\section{Why the general infinite-dimensional case is open}\label{sec:open}

\subsection{The three places where $\dim\cA<\infty$ is load-bearing}

\begin{enumerate}
\item \emph{The first projection (Lemma~6.1).} The induction starts from a nontrivial approximate projection, found as a fixed point of inversion on the approximate unitary group via the Lefschetz--Hopf theorem. That group is a compact manifold only when $\dim\cA<\infty$. In infinite dimensions the conclusion of the lemma is false even for genuine $C^*$ algebras: $C[0,1]$ and $C^*_r(F_2)$ have no nontrivial projections. So no argument can supply the first step; a proof would have to start from something else, for example an approximately commutative subalgebra close to some $C(K)$, which would need an approximate Gelfand theory that does not exist yet.
\item \emph{Termination (\S9).} Merging and extension stop because $\dim\cA$ is finite. Each step re-runs error reduction, which moves the whole map by $O(\eps)$, so an infinite sequence of steps need not converge.
\item \emph{Error reduction (Lemma~8.2).} The Newton step uses an exact contractive diagonal $D\in\cB\hat\otimes\cB$. A $C^*$ algebra with a bounded diagonal is finite-dimensional (contractible). For infinite-dimensional $\cB$ the natural replacement is Johnson's bounded approximate diagonal, which exists exactly when $\cB$ is nuclear, and which needs weak$^*$ limits in the target. For non-nuclear $\cB$ (such as $\Bo(\ell^2)$ or $C^*_r(F_2)$) there is no diagonal of any kind, and Ulam-type stability is known to fail in neighbouring problems. Note that Corollary~8.3 does allow an infinite-dimensional \emph{target}; what is missing is the construction of $\cB$.
\end{enumerate}

\subsection{No local obstruction}

\begin{proposition}\label{prop:local}
For every $k\in\mathbb{N}$ and $\delta>0$ there is $\eps(k,\delta)>0$ such that the following holds. For every $\eps(k,\delta)$-$C^*$ algebra $\cA$ and every $x_1,\dots,x_k$ in its unit ball, there are a $C^*$ algebra $\cB$ and $y_1,\dots,y_k\in\cB$ with $|\|w(x)\|-\|w(y)\||\le\delta$ for every $*$-word $w$ of length at most $k$, evaluated with a fixed bracketing.
\end{proposition}

\begin{proof}
If not, there are $k,\delta$ and $\eps_n$-$C^*$ algebras $\cA_n$ with $\eps_n\to0$ and tuples $x^{(n)}$ for which every pair $(\cB,y)$ fails. Fix a free ultrafilter $\mathcal{U}$. The ultraproduct $\prod_\mathcal{U}\cA_n$, with coordinatewise operations, is well defined because the products are uniformly bounded. It satisfies (2.1)--(2.6) with $\eps=\lim_\mathcal{U}\eps_n=0$, so it is a $C^*$ algebra. Take $\cB$ to be this ultraproduct and $y_i=[x_i^{(n)}]$. Then $\|w(y)\|=\lim_\mathcal{U}\|w(x^{(n)})\|$ for each of the finitely many words, so for $\mathcal{U}$-most $n$ all of them are within $\delta$, a contradiction.
\end{proof}

So any counterexample must be global: it can only be witnessed by configurations of unbounded size or complexity. By Theorem~\ref{thm:B} and the remark after it, it also cannot be assembled from finite-dimensional pieces in a bounded way.

\subsection{Where a counterexample would have to come from}

The phenomenon ``locally small error, globally far from exact'' is familiar from approximate group homomorphisms. Its degree-two form is the following.

\begin{proposition}[warm-up: quasimorphisms]\label{prop:qm}
Let $f\colon\Gamma\to\mathbb{R}$ be a homogeneous quasimorphism with defect $D$ that is not a homomorphism (such $f$ exist for $\Gamma=F_2$), and let $0<t<2\pi/D$. Then $\chi_t=e^{itf}$ satisfies $|\chi_t(gh)-\chi_t(g)\chi_t(h)|\le tD$. Yet for every Hilbert space $\mathcal{K}$ there is no homomorphism $\pi\colon\Gamma\to U(\mathcal{K})$ with $\sup_g\|\pi(g)-\chi_t(g)1\|<\sqrt3$.
\end{proposition}

\begin{proof}
Suppose $\pi$ exists. Let $\lambda=\chi_t(g)e^{i\alpha}\in\sigma(\pi(g))$. Then $\lambda^n\in\sigma(\pi(g^n))$ and $\chi_t(g^n)=\chi_t(g)^n$ by homogeneity, so $|e^{in\alpha}-1|<\sqrt3$ for all $n$. This forces $\alpha\in2\pi\mathbb{Z}$: otherwise some multiple of $\alpha$ lies in $[2\pi/3,4\pi/3]$ mod $2\pi$. Hence $\pi(g)=\chi_t(g)1$ for all $g$, so $\chi_t$ is multiplicative. Then $t(f(gh)-f(g)-f(h))\in2\pi\mathbb{Z}$, with absolute value at most $tD<2\pi$, so $f$ is a homomorphism, a contradiction.
\end{proof}

Associativity is a degree-three condition, so the analogue lives in $H^3_b$. Let $\Gamma$ be a group with a bounded real 3-cocycle $\beta$ whose class is nonzero in $H^3_b(\Gamma;\mathbb{R})$ but zero in $H^3(\Gamma;\mathbb{R})$. For closed surface groups $\Gamma=\pi_1\Sigma_g$ with $g\ge2$, $H^3(\Gamma;\mathbb{R})=0$ while $H^3_b(\Gamma;\mathbb{R})$ is infinite-dimensional (Soma). Write $\beta=dc$ with an unbounded 2-cochain $c$, and twist the group algebra by
\[
e_g\star_t e_h=e^{itc(g,h)}e_{gh},\qquad\text{so that}\qquad(e_g\star_te_h)\star_te_k=e^{-it\beta(g,h,k)}\,e_g\star_t(e_h\star_te_k).
\]
On basis vectors the associator is at most $t\|\beta\|_\infty$, which is arbitrarily small. One step of the analogue of Proposition~\ref{prop:qm} is elementary.

\begin{proposition}[graded rigidity]\label{prop:graded}
If $t\|\beta\|_\infty<\pi$, there is no $U(1)$-valued 2-cocycle $\sigma$ (that is, no associative twist) with $\sup_{g,h}|e^{itc(g,h)}-\sigma(g,h)|<2\sin(\pi/8)\approx0.77$.
\end{proposition}

\begin{proof}
Write $e^{itc}\sigma^{-1}=e^{i\gamma}$ with $|\gamma|<\pi/4$ and $\sigma=e^{i\tau}$. Then $tc=\tau+\gamma+2\pi m$ with $m$ integer-valued, and $d\tau$ is $2\pi\mathbb{Z}$-valued because $\sigma$ is a cocycle. Applying $d$ shows that $t\beta-d\gamma$ is $2\pi\mathbb{Z}$-valued. But $|t\beta-d\gamma|<\pi+4\cdot\pi/4=2\pi$, so $t\beta=d\gamma$. Hence $\beta=d(\gamma/t)$ with $\gamma/t$ bounded, i.e.\ $[\beta]=0$ in $H^3_b$, a contradiction.
\end{proof}

Two steps separate this from a counterexample to the infinite-dimensional theorem, and we could not complete either.
\begin{enumerate}
\item \emph{A $C^*$-type norm.} In $\ell^1(\Gamma)$ the basis-vector bounds extend to all elements, but $\ell^1$ violates the $C^*$ identity badly. For a norm satisfying (2.5) one needs the phases $e^{-it\beta}-1$ to be small as \emph{multipliers} on the relevant tensor products, which is much stronger than smallness in sup norm. Physically, this is Jackiw's non-associativity of magnetic translations in the presence of monopole density: in hyperbolic space, uniform charge density gives a uniformly bounded charge in every geodesic tetrahedron, so the associator is uniformly small, and the obstruction is a bounded volume-type class.
\item \emph{Ungraded rigidity.} Proposition~\ref{prop:graded} excludes nearby \emph{graded} associative twists. A counterexample needs to exclude every $C^*$ algebra, i.e.\ the analogue of the spectral step in Proposition~\ref{prop:qm}. For non-nuclear algebras there is no spectral theory of the relevant elements to borrow.
\end{enumerate}

\subsection{What we expect}

Our guess is that the infinite-dimensional statement is true for $\eps$-$C^*$ algebras that are ``nuclear-like'', where approximate diagonals are available. We suspect that it fails in general, with the counterexample living where bounded cohomology of non-amenable groups lives. This is a guess, not a result.

\section{What is verified in Lean}\label{sec:lean}

Both theorems are machine-checked in Lean~4 with Mathlib (toolchain \texttt{v4.35.0-rc4}). Each is stated as a single proposition \texttt{Claim} and proved as \texttt{main : Claim}, and each passes an audit. The audit scans for \texttt{sorry}, \texttt{admit} and new axioms, and checks that \texttt{main} proves exactly \texttt{Claim}. It confirms with \texttt{\#print axioms} that only Lean's standard axioms (\texttt{propext}, \texttt{Classical.choice}, \texttt{Quot.sound}) are used, and replays every module in an independent kernel checker. Results that are not formalized enter only as named hypotheses of \texttt{Claim}, never as axioms. The companion repository contains the Lean sources (folder \texttt{lean/}) and the file \texttt{LEAN\_CERTIFICATE.md}, which gives the formal statements, the assumptions and a step-by-step map from the text to the Lean lemmas.
\begin{itemize}
\item \emph{Theorem~\ref{thm:A}}, conditional on six named results. Three are textbook facts about completely positive maps: $\|\varphi\|_\cb\le\|\varphi(1)\|$ and $\varphi(x^\dagger)=\varphi(x)^\dagger$ for CP $\varphi$; Stinespring's theorem; and the existence of unitary 1-designs. The other three are Kitaev's Theorem~12.1, equation (12.6) and Theorem~10.5, each with its $O(\cdot)$ made explicit as a constant. Everything else is proved: the construction of $\Delta$, the properties (12.34)--(12.36) at every matrix level, Lemma~12.4, the choice of $\xi_j$, the construction of $\Upsilon$, both normalizations, and the case $\eta=0$. The constants come out explicit: $\eta_0=\min\bigl(\eta_1,\eta_2,\eps_3/(c_1+1),1/(2\cdot10^5\Lambda)\bigr)$ and $C=10^6\Lambda$, with $\Lambda=C_3c_1+c_2+1$, where $(c_1,\eta_1)$, $(c_2,\eta_2)$, $(C_3,\eps_3)$ are the constants of the three cited results. A non-vacuity check shows that the hypotheses hold on $\ell^2(\mathbb{N})$ for $\Phi(X)=\langle e_0,Xe_0\rangle1$, which is UCP, idempotent and of rank one.
\item \emph{Theorem~\ref{thm:B}}, its first sentence, conditional on Kitaev's Theorem~2.3. This includes Proposition~\ref{prop:sum}. A non-vacuity check shows that the hypotheses hold with infinitely many sectors, and that at $\eps=0$ the transcribed Definitions~2.1--2.2 rebuild Mathlib's \texttt{CStarAlgebra}.
\item Also formalized, with only the standard axioms but outside the certificate audit: the inverse of a $\delta$-isomorphism is a $\delta/(1-\delta)^3$-isomorphism (Kitaev's ``$\delta+O(\delta^2)$'' with an explicit constant), and Proposition~\ref{prop:graded} in full.
\end{itemize}
Not formalized: Lemma~\ref{lem:rank}, Corollary~\ref{cor:codes}, the second sentence of Theorem~\ref{thm:B}, the remark following Theorem~\ref{thm:A}, and Section~\ref{sec:open} apart from Proposition~\ref{prop:graded}. Kitaev's own proofs of Theorems~2.3 and~10.5 (Lefschetz--Hopf, H-spaces, the implicit function theorem on the approximate unitary group) are far beyond what Mathlib supports, so they enter only as hypotheses.

\end{document}
